\documentclass[book.tex]{subfiles}

\begin{document}

\chapter{Path Integrals}

\label{chap:path_integrals}
\noindent\rule{11cm}{0.4pt} \\
\textbf{Prerequisites:}  \autoref{chap:schrodinger} \\
\textbf{Difficulty Level:} ** \\
\noindent\rule{11cm}{0.4pt}
\vspace{5mm}

Path integral techniques provide a method of computing transition amplitudes for sophisticated quantum mechanical systems. Path integrals are useful in standard quantum mechanics, but invaluable in quantum field theory as we will learn in future chapters.

The propagator defines the amplitude of transition from quantum state $q$ at time $0$ to $q'$ at time $T$ according to Hamiltonian $\hat{H}$
\begin{equation}
K(q', T, q, 0) = \left \langle q' \middle |e^{-i\hat{H} T} \middle | q \right \rangle 
\end{equation}
The propagator in theory defines the complete time evolution of a quantum system. The challenge in practice though is that the operator exponential $e^{-i\hat{H}t}$ is likely very complicated and infeasible to compute. Let's expand the Taylor series of this operator
\begin{equation}
    e^{-i\hat{H}t} = 1 - i\hat{H}t - \hat{H}^2 t^2 + \cdots 
\end{equation}
If the time $t$ is small, we can use the linear Taylor expansion to compute the propagator as
\begin{equation}
    K(q', t, q, 0) \approx \langle q' | q \rangle - it \langle q' | H | q \rangle 
\end{equation}
Is there a way we can break down our more complex propagator $K$ using this linear Taylor expansion? To this, recall that if $\{|q_i\rangle \}$ is an orthonormal basis for Hilbert space $\mathcal{H}$ then
\begin{equation}
    \sum_i | q_i \rangle \langle q_i | = I
\end{equation}
Take our orthonormal basis to be the set of eigenstates of the position operator. Since there are uncountably many such operators, we will write this identity equation as
\begin{equation}
    \int dq | q \rangle \langle q | = I
\end{equation}
We can split apart the matrix exponential as follows (this is valid since $\hat{H}$ commutes with itself).
\begin{equation}
    e^{-i\hat{H}T} = e^{-i\hat{H}\frac{T}{2}} e^{-i\hat{H}\frac{T}{2}} 
\end{equation}
Note that using this identity, we can perform the calculation

%\textcolor{red}{TODO: Can this be made more rigorous?}
\begin{align}
    K(q', T, q, 0) &= \left \langle q' \middle | e^{-i\hat{H} T} \middle | q \right \rangle \\
    &= \left \langle q' \middle | e^{-i\hat{H} \frac{T}{2}} e^{-i\hat{H} \frac{T}{2}} \middle | q \right \rangle \\
    &=  \langle q'  | e^{-i\hat{H} \frac{T}{2}} \int dq_1 | q_1 \rangle \langle q_1 | e^{-i\hat{H} \frac{T}{2}}  | q  \rangle \\
    &= \int dq_1 \langle q'  | e^{-i\hat{H} \frac{T}{2}}  | q_1 \rangle \langle q_1 | e^{-i\hat{H} \frac{T}{2}}  | q  \rangle \\
    &= \int dq K\left (q', T, q_1, \frac{T}{2} \right ) K\left (q_1, \frac{T}{2}, q, 0 \right )
\end{align}
Now let $\delta = T/N$. You will show in the exercises 
\begin{equation}
    K(q',T, q, 0) = \int dq_1 \cdots dq_n K_{q', q_{n}} \cdots K_{q_1, q} 
    \label{pathexpansion}
\end{equation}
where $K_{q_i, q_{i-1}} = K(q_i, i\delta, q_{i-1}, (i-1)\delta)$

\begin{figure}[ht]
   % \includegraphics[width=\linewidth]{figures/Physics/quantum_mechanics/path_integrals/path_integral.png}
   \includegraphics[width=\linewidth]{figures/Physics/quantum_mechanics/path_integrals/Expansion of Paths (1).png}
   \caption{Diagrammatic Expansion of Paths.} %TODO: This figure is sourced from \url{https://cds.cern.ch/record/435548/files/0004090.pdf}. Replace with our own illustration.}
\label{fig:path_integral}
\end{figure}

You can intuitively think of this expansion as traveling over all possible paths from $q$ to $q'$. 


%\textcolor{red}{TODO: Introduce the concept of Wick rotation and discuss the schrodinger equation as the Wick rotation of the heat equation}

% TODO: Work out the types of random samples

\begin{lstlisting}[language=python]

sample_path : Rand[$\mathbb{R}^3 \to \mathbb{R}^3$]
path_prob: ($\mathbb{R}^3 \to \mathbb{R}^3$) $\to$ $\mathbb{R}$

def path_integral(q: $\mathbb{R}^3$, s: $\mathbb{R}^3$): $\mathbb{C}$
  amplitude = 0
  for i in range(N):
    path = sample_path(q, s)
    amplitude += path_prob(path)
  return amplitude
\end{lstlisting}

\subsection{Exercises}
\begin{enumerate}
    \item Prove that the eigenfunctions of the position operator $\hat{x}$ are given by the Dirac delta functions $\delta_{x_0} = \delta(x - x_0)$.
    \item Prove the expansion equation for paths given in Equation \ref{pathexpansion}.
\end{enumerate}


\section{}
\end{document}


